% Author record template for the reporting protocol, version 1.0.0
% Standalone usage: pdflatex author_record_template.tex
\documentclass[10pt]{article}

\usepackage[margin=20mm]{geometry}
\usepackage[T1]{fontenc}
\usepackage[utf8]{inputenc}
\usepackage{array}
\usepackage{booktabs}
\usepackage{longtable}
\usepackage{ragged2e}
\usepackage[hidelinks]{hyperref}

\newcolumntype{L}[1]{>{\RaggedRight\arraybackslash}p{#1}}
\setlength{\tabcolsep}{3pt}
\renewcommand{\arraystretch}{1.15}

\newcommand{\RecordHeader}{%
\toprule
\textbf{Step} & \textbf{Item} & \textbf{Answer} & \textbf{Entry} & \textbf{Location of details}\\
\midrule}
\newcommand{\Blank}{\emph{(fill in)}}

\begin{document}

\section*{Author record under the reporting protocol (version 1.0.0)}

\noindent\textbf{Paper title:} \Blank\\
\textbf{Version and date:} \Blank\\
\textbf{Corresponding author:} \Blank\\
\textbf{Code and data:} \emph{(URL, DOI, or repository tag/commit, or the access restriction)}

\medskip
\noindent\textbf{Instructions.} Complete one entry per item, starting with the claim, since it bounds every later entry. Each entry states the choice or value that was actually used, in one or two sentences, followed by the location of the full details; a pointer alone is not an entry, since it leaves the reader to reconstruct precisely the information that the record is meant to collect. Answer each item as \emph{value}, as \emph{not applicable} with the reason the component is absent from the workflow, or as \emph{not performed} with the consequence for the claim. Items marked \emph{(if applicable)} are answered as not applicable when the component is absent.

\setlength{\LTpre}{4pt}
\setlength{\LTpost}{8pt}

\begin{longtable}{@{}L{0.05\textwidth} L{0.20\textwidth} L{0.12\textwidth} L{0.36\textwidth} L{0.19\textwidth}@{}}
\RecordHeader
\endfirsthead
\RecordHeader
\endhead
\bottomrule
\endfoot
1 & Claim and evidence level & & & \\
2 & Stage ledger & & \emph{(see table below)} & \\
3 & Data & & & \\
3 & Splits & & & \\
4 & Model & & & \\
4 & Training & & & \\
4 & Generation & & & \\
5 & Evaluation & & & \\
5 & Feasibility & & & \\
6 & Baselines and ablations & & & \\
7 & Closed loop and uncertainty \emph{(if applicable)} & & & \\
7 & Reproducibility & & & \\
7 & Compute \emph{(if applicable)} & & & \\
7 & Governance \emph{(if applicable)} & & & \\
8 & Key numbers & & & \\
\end{longtable}

\subsection*{Stage ledger}
For iterative workflows, give one block of rows per cycle type and a campaign total.

\begin{longtable}{@{}L{0.16\textwidth} L{0.30\textwidth} L{0.08\textwidth} L{0.08\textwidth} L{0.14\textwidth} L{0.14\textwidth}@{}}
\toprule
\textbf{Stage} & \textbf{Rule or evaluator} & \textbf{In} & \textbf{Out} & \textbf{Evaluator calls} & \textbf{Cost per call}\\
\midrule
 & & & & & \\[1.2em]
 & & & & & \\[1.2em]
 & & & & & \\[1.2em]
 & & & & & \\[1.2em]
\bottomrule
\end{longtable}

\subsection*{Key-number summary for the main text}
\emph{One or two sentences: proposals generated; candidates and evaluator calls at each level of evidence; validated outcomes and the level at which they were validated.}

\end{document}
